\section{資料集配置與統計}
\label{sec:dataset}

為了全面評估模型在不同場景下的效能與泛化能力，本研究選用了三個具代表性的資料集系列，涵蓋了標準邏輯存取攻擊（ASVspoof）、真實世界場景（InTheWild）以及多語言環境（MLAAD）。各資料集的詳細統計數據，包含語言種類、生成系統數、語句數量及總時長等資訊，彙整於表 \ref{tab:dataset_stats}。

\begin{table}[H]
\centering
\caption{實驗資料集統計概覽}
\label{tab:dataset_stats}
\resizebox{\textwidth}{!}{%
\begin{tabular}{lcccccc}
\hline
\textbf{Dataset} & \textbf{Lang.} & \textbf{Systems} & \textbf{Dur.(hr)} & \textbf{Utterances} & \textbf{Spoof Ratio(\%)} & \textbf{Avg Dur.(s)} \\ \hline
ASVspoof19 LA & English & 19 & 109.7 & 121,461 & 89.72 & 3.25 \\
ASVspoof21 LA & English & 13 & 111.8 & 164,612 & 90.00 & 2.72 \\
ASVspoof21 DF & English & 100+ & 454.4 & 593,253 & 97.21 & 3.06 \\
In-The-Wild   & English & ? & 37.9 & 31,779 & 37.18 & 4.29 \\
M-AILABS      & 8 & 0 & 991.1 & 493,658 & 0.00 & 7.23 \\
MLAAD v2      & 23 & 52 & 160.2 & 72,000 & 100.00 & 7.79 \\
MLAAD v8      & 40 & 119 & 570.3 & 243,000 & 100.00 & 8.45 \\ \hline
\end{tabular}
} %
\begin{tablenotes}\small
\item \textit{註：}此處的「?」表示 In-The-Wild 資料集中偽造語音的生成系統數量無法精確統計：這些假樣本是從公開網路上蒐集之深偽影音剪輯彙整而成，原始生成模型並未一一標註，因此無法給出確切的 Systems 數量。\cite{muller2022does}
\end{tablenotes}
\end{table}

\begin{enumerate} \item \textbf{ASVspoof 系列資料集} \cite{wang2020asvspoof, yamagishi2021asvspoof}： 

% 作為本領域的標準基準，我們使用了 \textbf{ASVspoof 2019 Logical Access (LA)} 、\textbf{ASVspoof 2021 LA} 與 \textbf{DeepFake (DF)} 。該資料集均已預先劃分為訓練 (Train)、開發 (Dev) 與評估 (Eval) 子集。

ASVspoof (Automatic Speaker Verification Spoofing and Countermeasures Challenge) 為該領域公認的權威基準。本研究整合了 2019 年與 2021 年的任務場景，以評估模型在傳統攻擊與現代通訊變異下的表現：
    
    \begin{itemize}
        \item \textbf{ASVspoof 2019 Logical Access (LA)}：作為模型訓練的核心基準，該資料集模擬偽造語音透過通訊系統直接注入的情境。其語料源於 VCTK 資料庫，涵蓋由 17 種文字轉語音 (TTS) 與語音轉換 (VC) 演算法生成的樣本，包含神經聲學模型（如 WaveNet）生成的偽造特徵。
        
        \item \textbf{ASVspoof 2021 Logical Access (LA)}：延續 2019 LA 的語料基礎，但進一步引入真實電信網路（VoIP 與 PSTN）傳輸所產生的通道損耗。其經過 A-law、G.722 與 $\mu$-law 等七種不同編解碼器處理，旨在測試模型對於通訊系統所產生的假影（Artifacts）之強健性。
        
        \item \textbf{ASVspoof 2021 DeepFake (DF)}：專注於社群媒體與網路傳播場景，偽造樣本由超過 100 種欺騙演算法生成。此子集引入了大量有損壓縮（Lossy compression）處理，包含不同位元速率下的 mp3、m4a 與 ogg 格式，總時數約 454.4 小時，是目前對壓縮抗性要求最高的場景。
    \end{itemize}

\item \textbf{MLAAD (Multilingual Audio Anti-Spoofing Dataset)} \cite{muller2024mlaad}：
% 這是一個包含多語言的大型偽造語音資料集。由於 MLAAD 本身僅包含偽造語音 (Spoof only) 樣本，無法單獨用於訓練二元分類器。因此，在訓練階段，我們參考論文作者建議，將其與此資料集基於的上游全真音訊資料集 \textbf{M-AILABS} 進行混合，以構建完整訓練數據。而在測試階段，我們關注模型對該資料集中偽造語音的檢出率。

MLAAD 是一個多語言的大型偽造語音資料集，第二版涵蓋 23 種語言 (160.2 小時)，而第八版則擴增至 40 種語言 (570.3 小時)，所有音訊均以 22.05 kHz 格式輸出。其攻擊樣本生成基於 M-AILABS 語音庫中的真實語音，利用 52 種 (v2) 增加至 119 種 (v8) 先進文字轉語音（TTS）模型與 19 種不同架構（如 VITS、FastSpeech 等）製作，顯著提升了資料的多樣性，生成來源廣泛涵蓋 Coqui.ai 與 Hugging Face 等開源平台。由於 MLAAD 本身僅包含偽造語音 (Spoof only) 樣本，無法單獨用於訓練二元分類器。因此，在訓練階段，我們參考論文作者建議，將其與此資料集基於的上游全真音訊資料集 M-AILABS 進行混合，以構建完整訓練數據。而在測試階段，我們關注模型對該資料集中偽造語音的檢出率。

\item \textbf{InTheWild (ITW)} \cite{muller2022does}：
% 為了驗證模型在真實世界場景下的表現，本研究使用了 InTheWild 資料集。此資料集包含了從真實網路環境中收集的深偽音訊，具有多樣的背景噪音與壓縮失真，無預先定義的分割，本研究將其作為獨立的跨資料集測試使用。

為了評估模型在實驗室外的泛化能力，本實驗為此採用 InTheWild 資料集作為獨立的跨資料集測試使用。此資料集專門收集自公開的網路來源，用以彌補 ASVspoof 2019 僅基於 VCTK 錄音室語料庫所帶來的局限性。ITW 資料集總長 37.9 小時，精選了 58 位英語名人與政治人物的語音剪輯。其內容被劃分為 17.2 小時的偽造 (fake) 音訊與 20.7 小時的真實 (real) 音訊。偽造音訊是從 219 個公開可用的影片與音訊檔案中分割建構。所有音訊經標準化處理並轉換為 16 kHz 取樣率，平均片段長度約 4.3 秒。

\end{enumerate}

\section{模型配置與參數設定}
\label{sec:model_description}

\subsection{硬體環境與實驗控制} 
所有的實驗均在單張 \textbf{NVIDIA RTX 3090 (24GB VRAM)} 顯卡上進行。為了消除不同模型原始實作中因浮點數精度差異可能帶來的效能影響，並優化記憶體使用效率，本研究對實驗環境進行了以下兩項關鍵的標準化控制：

\begin{enumerate} \item \textbf{統一精度 (Unified Precision)}： 儘管 W2V-AASIST 與 MFA-Conformer 的原始實作採用 FP32 精度，本研究將所有實驗模型的運算精度統一設定為 \textbf{bfloat16}。此舉不僅確保了模型間比較的公平性，更大幅降低了大型模型 (如 SSL-AASIST) 的記憶體佔用，使得我們能夠在單卡環境下使用更大的 Batch Size 進行更穩定的訓練。

\item \textbf{統一損失函數 (Unified Loss Function)}：
為了應對訓練資料中真實與偽造語音可能存在的數量不平衡問題，我們摒棄了 SpeechPrompt 與 MFA-Conformer 原本使用的標準交叉熵 (CE)，將所有模型的損失函數統一固定為 \textbf{加權交叉熵 (Weighted Cross Entropy, WCE)}。這與 W2V-AASIST 的原始設定保持一致，確保所有模型在相同的優化目標下進行學習。

\end{enumerate}

\subsection{模型實作設定細節}

\textbf{SSL-AASIST (W2V-AASIST)}：該模型以原始波形 (Raw Waveform) 作為輸入（其深度學習防偽模型之發展脈絡請參見章節 \ref{ch:3-detection}）。由於其前端依賴參數高達 3 億的 wav2vec 2.0 XLSR 模型 \cite{conneau2020unsupervised}，為了適應單卡 3090 的記憶體限制並保持訓練穩定性，我們利用 bfloat16 帶來的空間優勢，將 Batch Size 從原始文獻 \cite{tak2022automatic} 的 14 提升至 40。根據線性縮放原則，我們相應地調整了學習率，將 Adam 優化器的 Learning Rate 從 $10^{-6}$ 調整至 $5\times 10^{-6}$。輸入長度 (Context Window) 固定為 4 秒。

\textbf{MFA-Conformer}：本研究復現了基於 Conformer 的架構，該模型以 \textbf{80 維梅爾頻譜圖 (Mel-spectrogram)} 作為輸入特徵（詳細文獻與架構回顧請見章節 \ref{ch:3-detection}）。實驗中使用 NVIDIA NeMo \cite{kuchaiev2019nemo} 提供的預訓練權重 \texttt{stt\_en\_conformer\_ctc\_small} 進行初始化，其餘超參數與參考文獻 \cite{wang2023pretraining} 保持一致：模型輸入長度設定為 5 秒，Batch Size 設定為 64。優化器採用 Adam，並搭配 Cosine Annealing 排程器 (Scheduler)，包含 Warmup 階段，最高學習率設定為 $10^{-3}$。

\textbf{SpeechPrompt v2}：SpeechPrompt 架構採用離散單元 (Discrete Tokens) 作為輸入（相關 Prompting 範式與 GSLM 架構請參見章節 \ref{ch:3-detection}）。由於離散化後的特徵序列長度大幅縮減，除 MLAAD 資料集中極少數超過 100 秒的異常樣本外，模型可將完整的語音長度納入 Context Window 中，無需進行截斷。實驗採用 \texttt{fairseq} 框架 \cite{ott2019fairseq} 進行訓練，由於該框架專為自然語言任務設定特性，會自動根據不同輸入長度動態調整 Batch Size 以盡可能貼近指定的 max token 數量，用以穩定訓練階段的 VRAM 占用，我們採用原始論文 \cite{chang2023speechprompt} 的參數進行訓練。
在 Prompt 設定方面，我們依據文獻建議 \cite{chang2023speechprompt}，啟用 \textbf{Deep Prompt} 機制，並將 Prompt Length 設定為 5。優化器學習率設定為 $5\times 10^{-3}$，由於 Prompt Tuning 過程的收斂速度極快，原始論文 \cite{chang2023speechprompt} 作法設定 Early Stopping 的 Patience 為 1。



\section{實驗結果}
\label{sec:model_performance}

\textbf{評估指標}：針對不同的測試情境，本研究採用兩種主要指標。首先，適用於 ASVspoof 與 InTheWild 等包含正負樣本的完整測試集，我們採用 \textbf{等錯誤率 (Equal Error Rate, EER)} 作為主要評估標準。EER 是指透過調整系統的決策閾值 ($\theta$)，當系統的錯誤接受率 (False Acceptance Rate, FAR) 等於錯誤拒絕率 (False Rejection Rate, FRR) 時的對應數值，其數學定義如下：
\begin{equation}
    EER = FAR(\theta) = FRR(\theta)
\end{equation}
\begin{equation}
    \text{where } FAR = \frac{\text{False Acceptances}}{\text{Total Impostor Attempts}}, \quad FRR = \frac{\text{False Rejections}}{\text{Total Genuine Attempts}}
\end{equation}
EER 的數值越低，代表系統的辨識與防偽性能越佳。其次，對於 MLAAD 資料集，因其全數為偽造語音（不包含真實語音），無法計算傳統的 EER。因此，我們改採 \textbf{偵測準確率 (Detection Accuracy, ACC)}，即計算模型成功將樣本判定為「Spoof」的比例作為準確率指標。

\begin{table}[H]
\centering
\caption{各模型在不同資料集與訓練設定下之效能表現}
\label{tab:main_results}
\begin{threeparttable}
\resizebox{\textwidth}{!}{%
\begin{tabular}{llcccc|cc}
\toprule
\multirow{2}{*}{\textbf{Model}} & \multirow{2}{*}{\textbf{Train Data \textbackslash\ Test on}} & \multicolumn{4}{c|}{\textbf{EER}} & \multicolumn{2}{c}{\textbf{Accuracy}} \\
\cmidrule(lr){3-6} \cmidrule(lr){7-8}
& & \textbf{ASV19 LA} & \textbf{ASV21 LA} & \textbf{ASV21 DF} & \textbf{ITW} & \textbf{MLAADv2} & \textbf{MLAADv8} \\
\midrule
\multirow{4}{*}{W2V-AASIST} 
& ASV19 & 0.27\% & 0.87\% & 5.57\% & 13.30\% & 83.24\% & 71.04\% \\
& ITW & 7.14\% & 10.90\% & 5.98\% & \cellcolor{gray!30}- & 55.48\% & 49.12\% \\
& MLAADv2 & 21.29\% & 30.26\% & 28.39\% & 18.77\% & \cellcolor{gray!30}- & \cellcolor{gray!30}- \\
& MLAADv8 & 26.85\% & 30.61\% & 23.48\% & 16.03\% & \cellcolor{gray!30}- & \cellcolor{gray!30}- \\
\midrule
\multirow{4}{*}{SpeechPrompt v2} 
& ASV19 & 18.50\% & 18.23\% & 28.94\% & 67.91\% & 99.65\% & 99.48\% \\
& ITW & 22.23\% & 25.10\% & 31.80\% & \cellcolor{gray!30}- & 98.26\% & 98.50\% \\
& MLAADv2 & 43.92\% & 44.68\% & 51.12\% & 63.31\% & \cellcolor{gray!30}- & \cellcolor{gray!30}- \\
& MLAADv8 & 45.10\% & 46.20\% & 50.80\% & 65.10\% & \cellcolor{gray!30}- & \cellcolor{gray!30}- \\
\midrule
\multirow{4}{*}{MFA-Conformer} 
& ASV19 & 1.12\% & 5.67\% & 16.00\% & 30.48\% & 92.61\% & 88.04\% \\
& ITW & 19.93\% & 26.48\% & 23.39\% & \cellcolor{gray!30}- & 37.31\% & 38.72\% \\
& MLAADv2 & 34.90\% & 41.34\% & 38.14\% & 31.34\% & \cellcolor{gray!30}- & \cellcolor{gray!30}- \\
& MLAADv8 & 45.27\% & 45.76\% & 42.98\% & 32.75\% & \cellcolor{gray!30}- & \cellcolor{gray!30}- \\
\bottomrule
\end{tabular}%
}
\begin{tablenotes}\small
\item \textit{註：}灰色儲存格 (\texttt{-}) 表示該測試集與訓練集相同，不進行評估。
\end{tablenotes}
\end{threeparttable}
\end{table}

本節將整理並對比各模型架構在不同測試集下的實驗結果，詳細數據請見表 \ref{tab:main_results}。本實驗之主要跨資料集測試設定借鑑自 MLAAD 研究文獻 \cite{muller2024mlaad}。前四個測試集（ASV19、ASV21 LA、ASV21 DF、ITW）以等錯誤率 (EER) 為評估指標；後兩個測試集（MLAADv2、MLAADv8）則以準確率 (Accuracy) 呈現。依據原始設計，我們僅對 \textbf{W2V-AASIST} 採用了其原始論文建議的資料增強 (DA) 策略。（在初步實驗中，我們亦嘗試對 MFA-Conformer 與 SpeechPrompt v2 施加相同的 DA 策略，但結果顯示這對模型效能並無顯著提升，因此為了保有其原始評估的純粹性，最終對比中這兩者皆維持不使用 DA 的設定）。

\textbf{實驗可重複性與真實效能探討}：在此必須說明，本研究所呈現之所有客觀評估數據，皆是基於三組不同隨機亂數種子（Random Seeds）進行獨立完整訓練後所取得之平均結果。本研究在測試階段不採取依賴單一最佳檢查點（Cherry-picking Checkpoints）的極端評估方式，因此，本實驗忠實呈現的結果與原文字面宣稱數據可能存在落差。例如：SpeechPrompt v2 原始文獻 \cite{chang2023speechprompt} 指出其模型於 ASV19 LA 能達到 13.5\% 的 EER，而本研究之平均結果為 18.50\%。即使我們完全依照原始文獻之設定採用 float16 精度，並針對驗證集頻繁挑選特定檢查點進行評估，亦僅能將 EER 推進至約 15\%，未能重現 13.5\% 的數據。此現象說明效能落差並非單純源自為了公平對比而採用的 bfloat16 統一精度配置。綜合上述，我們堅持採用更嚴格的三組平均評估標準，以避免選擇性保留極值的偏差，客觀地反映出模型面對未知攻擊時的真實防禦下限。

\textbf{基於標準資料集訓練之表現}：當以標準測試集 (ASV19) 進行訓練時，\textbf{W2V-AASIST} 在同領域及跨領域 (ASV21 LA/DF) 測試上均展現出顯著的優勢，EER 分別低至 0.27\%、0.87\% 與 5.57\%。\textbf{MFA-Conformer} 的表現居次，在 ASV19 上獲得 1.12\% 的 EER，但在跨領域的 DF 測試集上升至 16.00\%。相對而言，\textbf{SpeechPrompt v2} 在這幾個資料集上的表現明顯落後，EER 皆高於 18\%。而在真實場景測試 (InTheWild, ITW) 中，也觀察到類似的趨勢：以 ASV19 訓練的模型，在 ITW 上 W2V-AASIST 能維持 13.30\% 的 EER，而 MFA-Conformer 與 SpeechPrompt v2 則分別大幅退步達到 30.48\% 與 67.91\% 的錯誤率。

\textbf{基於跨領域資料集訓練之結果}：當改以 ITW 作為訓練資料時，所有模型在 ASV 系列上的表現皆普遍不如以 ASV19 訓練的成績。其中，W2V-AASIST 在 ASV19 LA 上維持了 7.14\% 的 EER，依然是三個架構中表現最好的；MFA-Conformer 與 SpeechPrompt v2 則分別為 19.93\% 與 22.23\%，顯示出基礎的反向跨資料集測試依然具備相當的挑戰性。此外，當分別以 MLAAD (v2 或 v8) 進行訓練後去測試其他資料集時，所有架構的 EER 都出現了巨幅的上升，效能全面下滑，顯示出個別模型在僅依賴單一資料集訓練時，其學習到的特徵極易侷限於該資料集的分佈特性，跨域泛化能力嚴重不足。

\textbf{針對 MLAAD 測試集之表現}：在 MLAAD (v2 與 v8) 的評估中，資料呈現出截然不同的反向趨勢。當同樣使用 ASV19 進行訓練時，在 ASV 與 ITW 系列表現最為墊底的 SpeechPrompt v2，卻在 MLAADv2 與 MLAADv8 上達到極高的準確率 (99.65\% 與 99.48\%)；MFA-Conformer 亦有約 88\%-92\% 的準確率；相較之下，在其他測試集表現最穩定的 W2V-AASIST 準確率反而僅約 71\%-83\%。相關現象的深層分析將保留至後續章節與其他消融實驗一起探討。

\section{消融實驗與延伸分析}
\label{sec:discussion}

基於前一節 \ref{sec:model_performance} 實驗結果所揭示的模型跨域效能衰退與在不同測試集上截然不同的表現趨勢，本節針對這些現象與模型架構特性進行深入的消融實驗與延伸分析。以下將依序探討：首先，比較 Prompt Tuning 與 Full Fine-Tuning 的效能極限，以釐清 SpeechPrompt 受限之主因（\ref{subsec:pt_vs_fft} 節）；其次，導入軟性離散特徵機制，探索彌補離散化特徵流失的可能性（\ref{subsec:soft_discrete} 節）；第三，提出可學習群集中心機制，進一步評估其對離散化特徵瓶頸的改善效果與潛在局限（\ref{subsec:learnable_centroids} 節）；第四，檢驗擴增預訓練模型規模對提升跨域偵測表現的實質影響（\ref{subsec:model_size_scaling} 節）；最後，評估將 MFA-Conformer 後端分類器升級為圖神經網路對多層次特徵融合的助益（\ref{subsec:backend_improvement} 節）。

\subsection{Prompt Tuning 與 Full Fine-Tuning 之極限對比}
\label{subsec:pt_vs_fft}
為了釐清 SpeechPrompt 性能受限的主因，究竟是源於 Prompt Tuning 方法本身的表達能力不足，還是作為骨幹網路的 GSLM 預訓練模型本身就不具備深偽偵測所需的聲學特徵，我們將測試基準擴展到全模型微調 (Full Fine-Tuning) 並將其視為在此特徵空間下的效能上限 (Upper Bound)。在 Full Fine-Tuning 的設置下，我們解除凍結 (Unfreeze) 整座 GSLM 語言模型，亦即除了 Verbalizer 外，原本被凍結的 12 層 Transformer Decoder 區塊的所有權重矩陣均開放於訓練階段進行參數更新。相關數據如表~\ref{tab:exp1_speechprompt} 所示。

\begin{table}[H]
\centering
\caption{SpeechPrompt v2 在 Prompt Tuning 與 Full Fine-Tuning (Upper Bound) 效能比較}
\label{tab:exp1_speechprompt}
\begin{tabular}{lcc}
\toprule
\textbf{Dataset} & \textbf{Prompt Tuning} & \textbf{Full Fine-Tuning (Upper Bound)} \\
\midrule
ASVspoof 19 LA & 18.50\% & 17.23\% \\
ASVspoof 21 LA & 18.23\% & 17.60\% \\
ASVspoof 21 DF & 28.94\% & 29.90\% \\
\bottomrule
\end{tabular}
\end{table}

實驗結果顯示，相較於僅使用 Prompt Tuning，全參數微調 (Full Fine-Tuning) 僅在 LA 任務上帶來小幅度的進步（EER 由 18.50\% 降至 17.23\%），而在深偽難度較高的 DF 任務上甚至出現略微退步（EER 由 28.94\% 上升至 29.90\%）。這項結果印證了兩個關鍵洞察：(1) \textbf{Prompting 的參數效率極高}：僅訓練極少量的 Prompt 參數即可達到接近全系統微調的效能水準，並且在面對未知的跨域攻擊時，能藉由避免極度過擬合 (Overfitting) 來維持同等的防禦下限；(2) \textbf{瓶頸在於前端特徵}：即使開放所有基礎語言模型參數進行微調，也無法取得跨越式的大幅效能進展。這暗示著效能天花板並非源於後端語言模型的學習能力不足，而是前端 S2u (Speech-to-unit) 階段的離散化轉換過程造成了防偽所需的細微連續聲學特徵流失，使得模型在先天上即缺乏辨別真假的關鍵輸入訊號。

儘管 SpeechPrompt 在極限偵測準確度上受限於離散特徵天花板，未能超越 W2V-AASIST，但其在「參數效率 (Parameter-Efficient)」上展現的壓倒性優勢，正是未來實務部署的重要指引。
在可訓練參數量方面，相較於 W2V-AASIST 必須對前端 Wav2Vec 2.0 (XLSR) 包含數十層卷積與 Transformer 區塊進行全參數微調（高達 \textbf{317M}），以及 MFA-Conformer (Small) 包含 16 層 Conformer 模塊需訓練約 \textbf{14.6M} 個參數，SpeechPrompt 採用 Prompt Tuning 機制大幅降低了訓練成本。它凍結了龐大的 GSLM 骨幹，僅需訓練 Input Prompts、各層 Transformer 的 Deep Prompts 附加向量與最後的 Verbalizer 線性層。在 Prompt Length $X=5$ 的設定下，其可訓練參數僅為 \textbf{129,640 (約 0.13M)}。這意味著 SpeechPrompt 的參數量僅為 W2V-AASIST 的 \textbf{0.04\%}，極大程度地降低了模型儲存與部署的門檻。

\subsection{軟性離散特徵輸入架構評估}
\label{subsec:soft_discrete}
為了驗證「離散化導致特徵流失」的假設，並嘗試在不改變模型主體架構的前提下找回遺失的聲學細節，我們提出了一種軟性離散特徵 (Soft Discrete Features) 的輸入機制。在原始的 SpeechPrompt 中，HuBERT 輸出的連續特徵會被硬性分配給最接近的矩陣中心點（Hard Quantization）。在本實驗中，我們改為計算特徵與各個中心點的歐氏距離，並透過 Softmax 函數將其轉換為相似度權重，進而將各個中心點特徵進行加權求和，作為下游語言模型的平滑輸入。

給定時間點 $t$ 的連續聲學特徵 $z_t \in \mathbb{R}^D$ 以及 $K$ 個預訓練的群集中心 $\{c_i\}_{i=1}^K$，其對應於第 $i$ 個群集中心的權重 $p_{t,i}$ 定義為：
\begin{equation}
    p_{t,i} = \frac{\exp(-\|z_t - c_i\|^2 / \tau)}{\sum_{j=1}^K \exp(-\|z_t - c_j\|^2 / \tau)}
\end{equation}
其中，$\tau$ 為控制分佈平滑程度的溫度參數 (Temperature)。最終的軟性特徵表示 $\tilde{z}_t$ 等於各中心點之加權線性組合：
\begin{equation}
    \tilde{z}_t = \sum_{i=1}^K p_{t,i} c_i
\end{equation}
我們預期透過此一平滑映射方式，能在一定程度上保留原始音訊在特徵空間中的相對位置資訊，從而補償硬性量化所帶來的資訊損失。

\begin{table}[H]
\centering
\caption{軟性離散特徵輸入機制之效能 (EER) 比較}
\label{tab:exp2_soft_discrete}
\begin{tabular}{lccc}
\toprule
\textbf{Model Config} & \textbf{ASV19 LA} & \textbf{ASV21 LA} & \textbf{ASV21 DF} \\
\midrule
SpeechPrompt Baseline & 18.50\% & 18.23\% & 28.94\% \\
Soft Discrete Baseline & 15.90\% & 17.34\% & 25.95\% \\
\bottomrule
\end{tabular}
\end{table}

實驗設定採用 Softmax 溫度參數 $\tau = 2$ ，學習率 (Learning Rate) 為 $1\times 10^{-3}$ 進行微調。實驗結果如表~\ref{tab:exp2_soft_discrete} 所示，此架構在 ASV19 LA 資料集上取得了 15.90\% 的等錯誤率 (EER)。雖然相較於原始硬性量化架構的 18.50\% 有所改善，且在其他跨域資料集上的表現 (ASV21 LA: 17.34\%, ASV21 DF: 25.95\%) 也優於硬性量化，但仍受限於靜態的群集中心佈局，未能完全突破前端特徵離散化所造成的效能瓶頸。

\subsection{可學習群集中心對離散化特徵瓶頸之改善}
\label{subsec:learnable_centroids}
基於上一節軟性離散特徵僅達到 15.9\% EER 的實驗結果，其效能瓶頸的根源在於群集中心本身仍為預訓練階段以 K-Means 演算法所產生的固定陣列，無法因應下游深偽偵測任務的需求而調整。為此，我們進一步提出兩種可學習的群集中心機制，使模型能端到端地 (End-to-End) 動態優化特徵投影空間，從而更有效地捕獲區分真偽語音的關鍵連續特徵。兩種方法皆以 HuBERT 輸出的 768 維連續特徵張量序列作為輸入，取代原始的離散 Token 序列。

\begin{table}[H]
\centering
\caption{可學習群集中心機制之效能 (EER) 比較}
\label{tab:exp3_learnable_centroids}
\begin{tabular}{lccc}
\toprule
\textbf{Model Config} & \textbf{ASV19 LA} & \textbf{ASV21 LA} & \textbf{ASV21 DF} \\
\midrule
SpeechPrompt Baseline（硬性量化） & 18.50\% & 18.23\% & 28.94\% \\
\midrule
(1) 群集中心殘差微調 (Centroid Residual Tuning) & 5.53\% & 12.88\% & 16.67\% \\
(2) 直接線性映射 (Direct Linear Mapping) & 3.84\% & 9.08\% & 14.44\% \\
\bottomrule
\end{tabular}
\end{table}

\textbf{(1) 群集中心殘差微調 (Centroid Residual Tuning)}：此方法保留預訓練階段建立的 K-Means 離散字典 $C_{fixed}$，但在其上疊加一個可學習的殘差參數矩陣 $\Delta C$，使訓練過程中的有效群集中心為 $C_{tuned} = C_{fixed} + \Delta C$。透過僅學習殘差項，模型得以在保留原始語音單元先驗知識的基礎上，針對深偽偵測任務微調中心點位置。實驗設定學習率為 $1\times 10^{-4}$，Softmax 溫度參數為 1.0。多組亂數種子實驗結果如表~\ref{tab:exp3_learnable_centroids} 所示，此方法成功突破了靜態群集機制所造成的效能底線，在未引發災難性遺忘 (Catastrophic Forgetting) 的情況下，將 ASV19 LA 的 EER 由軟性離散的 15.90\% 大幅降低至 \textbf{5.53\%}，其餘測試集亦有顯著進步。

\textbf{(2) 捨棄量化器之直接線性映射 (Direct Linear Mapping without Quantizer)}：此方法的核心概念在於完全移除原本擔任特徵離散化角色的 K-Means 量化器 (Quantizer)，並捨棄基於 L2 歐氏距離的群集配對機制。取而代之的是，我們在模型前端直接引入一個不受約束的線性神經網路層 (例如：\texttt{nn.Linear(768, 100)})，將 HuBERT 輸出的 768 維連續特徵直接降維映射至 100 維的特徵空間。此設計完全繞過了離散化瓶頸，賦予模型最大的自由度來學習深偽偵測任務導向的特徵轉換。如表~\ref{tab:exp3_learnable_centroids} 所示，在此架構下 ASV19 LA 的 EER 進一步降低至 \textbf{3.84\%}，相較於原始包含量化器的離散化架構 (18.50\%) 實現了巨大的效能躍升。

上述兩組實驗雖然從量化指標上展示了顯著的改善，但也暴露了數項值得審慎考量的潛在風險。首先，\textbf{過擬合風險}：捨棄量化器之直接線性映射方法完全解除了群集空間的結構約束，使得模型容易過度適配 ASV19 LA 訓練集中特定的偽造演算法特性與錄音條件。在跨域測試集驗證 (如 ASV21 DF 取得 14.44\% EER) 中，其表現是否能持續穩定超越保留先驗知識之模型仍有待商榷。其次，\textbf{語音先驗知識的喪失}：原始的 K-Means 群集中心係由 LibriSpeech 960 小時未標註語料透過 HuBERT 自監督學習所建立，蘊含了豐富的語音單元結構資訊。直接線性映射方法完全繞過了此先驗知識，在訓練資料量有限的情境下，模型可能更傾向於學習到與深偽偵測無關的統計捷徑，而非真正的偽造特徵。

綜合而言，由軟性離散特徵 (15.90\%)、群集中心殘差微調 (5.53\%) 到捨棄量化器之直接線性映射 (3.84\%) 的遞進式效能提升，清楚地呈現出一條趨勢：隨著 HuBERT 連續特徵至語言模型輸入之間的映射方式從僵化的硬性量化逐步放寬為可學習的投影，下游語言模型所能獲取的有效聲學資訊量持續增加，偵測效能亦隨之顯著改善。此結果有力地印證了先前的假設——硬性離散化過程中的量化資訊損失，正是制約 SpeechPrompt 架構偵測能力的核心瓶頸。然而，放寬約束所帶來的效能提升與上述過擬合及先驗知識喪失之風險之間如何取得平衡，仍是後續研究需面對的關鍵議題。

\subsection{預訓練模型擴容對跨域偵測之影響}
\label{subsec:model_size_scaling}
MFA-Conformer (Small) 在 ASV21 DeepFake (DF) 任務上與 W2V-AASIST 存在顯著落差（EER 13.43\% vs 5.57\%）。為了驗證此差距是否源於模型容量不足，我們將前端替換為更大型的 \texttt{Conformer-Medium} 與 \texttt{Conformer-Large}。實驗評估數據如表~\ref{tab:exp3_model_size} 所示：即便將模型規模大幅提升，在 DF 測試集上的 EER 並未出現顯著下降，性能呈現飽和狀態（甚至 \texttt{Medium} 的錯誤率還略高於 \texttt{Small}）。這暗示了單純增加神經網路參數量並無法解決該架構在跨域深偽偵測上的關鍵瓶頸。

\begin{table}[H]
\centering
\caption{MFA-Conformer 不同預訓練模型規模之參數與效能 (EER) 比較}
\label{tab:exp3_model_size}
\resizebox{\textwidth}{!}{%
\begin{tabular}{lcccccc}
\toprule
\textbf{Model Size} & \textbf{Params} & \textbf{Layers} & \textbf{Dim.} & \textbf{ASV19 LA} & \textbf{ASV21 LA} & \textbf{ASV21 DF} \\
\midrule
Conformer-Small & 13.1M & 16 & 176 & 1.25\% & 4.58\% & 13.43\% \\
Conformer-Medium & 30.7M & 18 & 256 & 1.13\% & 3.84\% & 14.20\% \\
Conformer-Large & 121.5M & 18 & 512 & 1.80\% & 6.52\% & 13.31\% \\
\bottomrule
\end{tabular}%
}
\end{table}


\subsection{前端多層次特徵與圖神經網路後端之融合}
\label{subsec:backend_improvement}
我們推測 MFA 原始架構中簡單的統計池化與線性分類器可能無法完全整合 Conformer 抽取的複雜特徵，因此嘗試將後端替換為 AASIST 架構。訓練配置方面，我們與基線模型保持完全一致，包含相同的優化器、學習率、訓練週期數與資料增強設定，僅將後端分類器從線性分類器替換為 AASIST 圖神經網路，以確保實驗結果的差異僅來自於後端架構的變更。在實驗過程中，我們對於特徵選取與丟棄策略進行了多組不同的測試對比，分別評估了僅使用最後一層、最後兩層、最後三層與全部層次特徵之效果，最終發現將 MFA-Conformer 中各層的特徵 (All layer) 全數導入後端能得到最佳的效能提升。基於此「All layer」設定的實驗數據如表~\ref{tab:exp4_mfaconformer} 與原始基礎對齊呈現。

\begin{table}[H]
\centering
\caption{MFA-Conformer 不同後端分類器架構之效能 (EER) 比較}
\label{tab:exp4_mfaconformer}
\begin{tabular}{lccc}
\toprule
\textbf{Model Backend} & \textbf{ASV19 LA} & \textbf{ASV21 LA} & \textbf{ASV21 DF} \\
\midrule
Linear Classifier (Baseline) & 1.12\% & 5.67\% & 16.00\% \\
AASIST Architecture (using All layer) & 0.86\% & 3.94\% & 12.71\% \\
\bottomrule
\end{tabular}
\end{table}

結果顯示這項改良帶來了顯著的性能提升：ASV19 LA 的 EER 從 1.12\% 下降至 0.86\%；ASV21 LA 亦由 5.67\% 改善至 3.94\%；在處理難度最高的 ASV21 DF 資料集時，表現進步到 12.71\%。這證明原先簡易的分類網絡限制了 MFA-Conformer 的潛能，透過引入圖神經網路等更強大的特徵融合機制（如 AASIST），確實能更有效率地聚合多層次的前端特徵，並進一步增強防偽系統的穩健性。